% chapters/55_llm_agent_glossary.tex
% Part of: AI / LLM / Agents / Hardware Notes

\section{General-purpose LLM + Agent Map}

\subsection{How to use this chapter}

This chapter is a navigation layer for the rest of these notes: quick definitions, how concepts connect, and where each idea belongs (training, inference, or agent systems).

\subsection{Where each concept belongs in this repo}

\textbf{LLM training and adaptation}
\begin{itemize}[leftmargin=1.5em]
    \item \textbf{Pretraining} (Chapter 11): learning broad language/world patterns from large raw corpora.
    \item \textbf{Fine-tuning, SFT, LoRA} (Chapter 12): specialization for domain/task behavior.
    \item \textbf{RLHF / preference optimization} (Chapter 13): shaping behavior with preference signals.
    \item \textbf{Distillation} (cross-cutting, mostly training): transferring behavior from teacher to student.
\end{itemize}

\textbf{Inference and deployment}
\begin{itemize}[leftmargin=1.5em]
    \item \textbf{Inference and decoding} (Chapter 15): temperature, top-k, top-p, caching, batching.
    \item \textbf{Quantization} (Chapter 16): reducing numerical precision for speed and memory.
    \item \textbf{RAG} (Chapter 18): retrieval pipeline to ground generation at inference time.
\end{itemize}

\textbf{Agent systems}
\begin{itemize}[leftmargin=1.5em]
    \item \textbf{Agent fundamentals} (Chapter 19): loop, planning, acting, reflection.
    \item \textbf{Tool use / function calling / MCP} (Chapter 20): external actions and integrations.
    \item \textbf{Memory and planning} (Chapters 22--23): state, persistence, and decision flow.
\end{itemize}

\subsection{Distillation vs Quantization (high-value distinction)}

\begin{itemize}[leftmargin=1.5em]
    \item \textbf{Distillation}: changes what the model learns by training a smaller student to imitate a larger teacher.
    \item \textbf{Quantization}: changes how the learned weights are stored/executed by using lower precision.
\end{itemize}

Rule of thumb: distillation is mostly a \textit{training-time} method; quantization is mostly an \textit{inference/deployment-time} method.

\subsection{Distillation: practical trade-offs}

\textbf{Pros}
\begin{itemize}[leftmargin=1.5em]
    \item Better cost-latency-quality trade-off than training small models from scratch.
    \item Good path to task-specific students (coding, math, support workflows).
    \item Cleaner teacher outputs can provide a strong learning signal.
\end{itemize}

\textbf{Cons}
\begin{itemize}[leftmargin=1.5em]
    \item Teacher errors and style artifacts transfer to the student.
    \item Student capability ceiling is often tied to teacher quality.
    \item Overusing synthetic outputs can reduce diversity and robustness.
\end{itemize}

\subsection{Quantization: practical intuition + numeric example}

Quantization compresses weight values into fewer buckets.

\begin{notebox}[title=Bucket example]
Suppose a layer's weights in one tensor range from $-1.0$ to $1.0$, and we use 4-bit quantization.

4-bit gives $2^4 = 16$ buckets.

Bucket width is:
\[
\Delta = \frac{1.0 - (-1.0)}{16 - 1} = \frac{2.0}{15} \approx 0.1333
\]

For a weight $w = 0.18372642$:
\[
q = \text{round}\left(\frac{w - (-1.0)}{\Delta}\right)
  = \text{round}\left(\frac{1.18372642}{0.1333}\right)
  = 9
\]

Dequantized value:
\[
\hat{w} = -1.0 + 9 \cdot 0.1333 \approx 0.1997
\]

So the exact value $0.18372642$ is stored as bucket index $9$, then reconstructed approximately as $0.1997$.
\end{notebox}

\textbf{Why this helps}
\begin{itemize}[leftmargin=1.5em]
    \item Lower memory footprint (VRAM/RAM)
    \item Faster inference on many hardware targets
    \item Enables larger models on local/edge devices
\end{itemize}

\textbf{What can degrade}
\begin{itemize}[leftmargin=1.5em]
    \item Multi-step reasoning, coding precision, and long-context stability can drop at aggressive bit-widths.
\end{itemize}

\subsection{Minimal glossary (highest-yield)}

\begin{description}[leftmargin=1.5em, style=nextline]
    \item[Pretraining] Learn broad priors from large raw corpora.
    \item[Fine-tuning] Specialize a pretrained model.
    \item[Distillation] Teach a smaller model from a stronger model.
    \item[Quantization] Compress numeric representation for efficiency.
    \item[RAG] Retrieve documents first, then generate grounded answers.
    \item[Agent] A model-driven system that plans and uses tools.
\end{description}

\subsection{Agent types taxonomy}

Five canonical agent types by increasing capability (Russell \& Norvig):

\begin{description}[leftmargin=1.5em, style=nextline]
  \item[Simple reflex agent]
    Condition-action rules only. No memory, no world model. Works solely in fully
    observable deterministic environments. Example: thermostat, traffic light.

  \item[Model-based reflex agent]
    Maintains an internal world model updated each turn; handles partial observability
    by tracking unobserved state. Still uses condition-action rules for decisions.

  \item[Goal-based agent]
    Has explicit goals; uses search or planning to find action sequences that achieve
    them. Requires reasoning about future states. Example: route-finding robot.

  \item[Utility-based agent]
    Extends goal-based agents with a utility function (real-valued desirability per
    state). Enables rational choice when goals conflict or are achievable to varying
    degrees.

  \item[Learning agent]
    Improves performance from experience. Components: performance element (acts),
    critic (evaluates), learning element (updates), problem generator (explores).
    LLMs trained via RLHF/DPO are learning agents at training time; at inference time
    they are typically goal- or utility-based.
\end{description}

\subsection{Agent architecture components}

A complete agent system integrates five functional subsystems:

\begin{center}
\begin{tikzpicture}[
  comp/.style={rectangle, rounded corners=4pt, draw=aiblue, fill=infobg,
               minimum width=2.6cm, minimum height=0.8cm, align=center,
               font=\small\bfseries},
  mem/.style={rectangle, rounded corners=4pt, draw=aipurple, fill=codebg,
              minimum width=2.6cm, minimum height=0.75cm, align=center,
              font=\small\bfseries},
  arr/.style={-{Stealth[length=6pt]}, thick, color=aigray},
  node distance=0.8cm and 1.4cm
]
  \node[comp] (percept) {Perception};
  \node[comp, right=1.3cm of percept] (reason)  {Reasoning};
  \node[comp, right=1.3cm of reason]  (plan)    {Planning};
  \node[comp, right=1.3cm of plan]    (exec)    {Execution};
  \node[mem,  below=1.0cm of reason, xshift=1.1cm] (mem) {Memory};

  \draw[arr] (percept) -- (reason);
  \draw[arr] (reason)  -- (plan);
  \draw[arr] (plan)    -- (exec);
  \draw[arr] (reason.south)  -- (mem.north west);
  \draw[arr] (plan.south)    -- (mem.north east);
  \draw[arr] (mem.west) -- ++(-0.5,0) |- (reason.south west);
\end{tikzpicture}
\end{center}

\begin{itemize}[leftmargin=1.5em]
  \item \textbf{Perception:} visual (CV), auditory (ASR), textual (NLP/NER),
    environmental (sensor fusion), and predictive (anticipating future states).
  \item \textbf{Reasoning:} context understanding, pattern recognition --- the LLM
    forward pass, optionally with chain-of-thought (Chapter~17).
  \item \textbf{Planning:} goal-setting, decision-making. ReAct, ReWOO, MCTS (Chapter~23).
  \item \textbf{Execution:} API calls, tool invocations, code execution (Chapter~20).
  \item \textbf{Memory:} working (context window), episodic (logs), semantic (vector DB),
    procedural (RL / fine-tuning) (Chapter~22).
\end{itemize}

\subsection{One mental model}

\begin{itemize}[leftmargin=1.5em]
    \item Pretraining gives the base capability.
    \item Fine-tuning aligns it to a job.
    \item Distillation lowers serving cost.
    \item Quantization lowers runtime footprint.
    \item RAG adds external knowledge at inference time.
    \item Agent tooling turns answers into actions.
\end{itemize}
